\section{The Fire of Creation}

\begin{quotex}
This review by \textbf{Julius Evola} of J J van de Leeuw's book \textbf{The Fire of Creation}, in its Italian translation, originally appeared in \emph{Bilychnis}, volume XXXII, August-September, 1928.

\textbf{J J van de Leeuw} was a Theosophist and a priest in the Liberal Catholic Church\footnote{\url{http://www.thelccusa.org/}}. In van de Leeuw's understanding of the Trinity, the Father is the Will, the Son the Logos, and the Holy Spirit is Thought. Evola sees in this book an echo of Joachim of Fiore\footnote{\url{https://en.wikipedia.org/wiki/Joachim_of_Fiore}}, a saintly Abbott and highly regarded by \textbf{Dante}. In this review, Evola picks out one of the more interesting themes of the book, summarizing it rather well. Since the Liberal Catholic Church, just like the Roman, claims apostolic succession, we could read this as what Evola wished the Catholic Church would teach. Consider it his contribution to Vatican II, had he been allowed.
\end{quotex}


Though belonging to the stream of Blavatsky's anglo-indian theosophy, about which we have written not very favorably in the past, this work, whose Italian translation we draw attention to, represents something deep, clear, and well sketched out. The usual mythological cosmologies are not found in it, nor does he stray into clairvoyance, reincarnation, or similar topics. Van de Leeuw instead proposes a \emph{practical} point of view, of direct spiritual experience.

The mystical prophecies of \textbf{Joachim of Fiore} had earlier announced: ``The Kingdom of the Father has passed, that of the Son is passing, the Kingdom of the Holy Spirit is on the point of arising.'' And Van de Leeuw declares precisely that today we are truly at the dawn of the coming of the third person in the dialectic of the divine; the direction of creative activity in motion, of the universe in which \emph{dynamis} or power is the fundamental note, and the immanent presence of the ``cosmic Fire'' is characteristic of this. The evolutionist and becomingistic direction, the meaning of liberty, of independence from every creed, of individual initiative; the more and more prevalent tendency to move from action to knowledge rather than from knowledge to action until now, are, in modern times, so many signs of it. Those who have been touched by the ``Fire of Creation'' throw themselves into action, confident of drawing their certainty and their religion from it: they no longer say ``we pray'', but they say instead ``we work''. The author insists greatly on this aspect of
\emph{experience}, and not of philosophy or belief. One is put in relationship as with a conductor charged with a current of spiritual life, from which proceeds a discharge that dynamizes all our being. In the experience of the ``Fire of Creation'' one has the sense of a superior energy of something previously felt, and we are electrified through action. In such a moment we feel not only that we desire to do some things, but that we \emph{can} do them: it seems to us that no obstacle can ever resist this enormous energy that we now feel inside of us.

The author seeks moreover to justify in terms of inner experience the same concept of the ``Divine Mother'': it would be an aspect complementary to the Fire or to the Holy Spirit, to which it stands in the same relationship of the feminine to the masculine. The sense of hot activity producing its happening in something, the perception of being the same nature and of transforming within ourselves the radiant powers, the creative activity of the creative Fire in the fruitfulness of advancement and in abundance of beauty and form. Close to the impulse to act, the direction that we can make of the exultant joy that is awakened when we know the ``Holy Spirit'' principle, the complementary aspect ``Mother'' is instead a process of intimate transformation that renders productive what used to be dead.

The author also gives examples of techniques to produce these experiences. The presupposition is that the ordinary world is only a projected image of our consciousness through the mind and senses. We confuse this image with reality. It is then a question of controlling such an idiosyncrasy of our human manufacture, that makes us project the consciousness of what is inside us around ourselves into the world, and to reconverge at the center, instead of losing ourselves in the contemplation of our image of the world, just as the prisoners of Plato's cave stare at the shadows on the walls. It is possible to concentrate our Self into a very small point of consciousness and to pass through that point into the world of the Real, of the ``things that are'', and no longer our exteriorized images of it. We have to push ourselves ahead, and having withdrawn our consciousness from its image of the world and stopped its transcendental imaginative faculty, reintegrating ourselves in that absolute immanence, from which they originate, at the same time, the sense of the Fire of Creation and the vision of the Real in the world.


\flrightit{Posted on 2013-04-10 by Cologero}

\begin{center}* * *\end{center}

\begin{footnotesize}
\begin{sffamily}
\texttt{Jason-Adam on 2013-04-11 at 13:27 said:}

I have not read this book but desire to after reading Evola's positive review.

Let me add though that the Liberal Catholic Church is a heretical breakaway sect based on Free-Masonry and rebellion against Vatican I and is not an apostolic Church.

\hfill

\texttt{Cologero on 2013-04-11 at 16:08 said:}

I suspect, Jason-Adam, that Evola's review is better and more inspiring than the book itself, which may disappoint. Unfortunately, I can't locate my copy of the book which I had read quite some time ago.

\hfill

\texttt{Jason-Adam on 2013-04-11 at 17:35 said:}

This is tangental but I feel it necessary to include this on a page devoted to Evola and the Church -- I hope Avery will see this and comment as well :

Avery made an interesting remark that all the political Traditionalists are anti-Christian whereas the Christian Traditionalists are not involved in political movements. I believe this is because we are seeing the difference between brahmans (the Christians) and kshatriyas (the pagans). I offer as an opinion the reason why modern day kshatriyas in the west are driven to despise the Church is because the Church has failed to maintain the path of action, the old time warrior Christianity of the Crusades. What Catholic of today would praise Godfrey de Bouillon the saviour of Jerusalem ? One of my favourite works of epic poetry is Gerusalemme Liberatta of Tasso's, would modern Catholics see themselves in that work ? Would any respectable theologian say with St Bernard that it is good and right to kill the enemies of the Catholic religion ? The general effeminancy and passivity of the modernist Church is turning away men of fighting spirit.

It is ironic that the main reason I am so drawn to catholicism, true catholicism, is that compared to the contemplative spirit of the eastern church, western catholicism is a religion of action as Soloviev eloquently explained in Russian and the Universal Church.

\hfill

\texttt{Michael on 2013-04-11 at 20:48 said:}

Van de Leeuw's teaching sounds attractive. Who would not want to be filled with the power to overcome all obstacles? But chasing after this kind of experience is probably a waste of time. Better to just be men of action in fact rather than waiting for some move of the Spirit, and we will find that it has been the Spirit moving us all along after all.

\hfill

\texttt{Cologero on 2013-04-12 at 08:34 said:}

Yes, Jason-Adam, you have identified the cause of the break up. I've noticed, however, that even the neo-pagans display images of Christian knights on their web sites, so they may not be as far from you as you may think. This split is bad for both sides. The ``brahmans'', as you call them even though they have not preserved the entire spiritual heritage, are attracting spirits of a certain type. Occasionally I check their writings and what is remarkable is the vitriol and antipathy they display toward their fellow religionists who doesn't accept that the church began --- or recreated itself anew --- in 1965. They are really modernists, perhaps with little quirks such as opposition to abortion; otherwise, spiritually they are quite similar.

The ``kshatriyas'', in their position, give up 1500 years of their history along with the deep metaphysical, spiritual, and moral teachings, etc. It seems to me that the battle needs to move to a different front and the spirituality appropriate for the kshatriya given its proper place.

\hfill

\texttt{Avery Morrow on 2013-04-12 at 12:57 said:}

I agree that the split is bad for both sides. I will give Jason-Adam's question a shot. For me, Hannah Arendt's \emph{The Human Condition} has been a key to this pagan-Christian puzzle. To summarize just some of her work: the ancients considered labor an animal necessity, and only who had absolutely no need for work could be considered freemen and citizens, who were suitably free of private hindrances to devote their lives to public work. So far, so Evola. Christianity cultivated this disgust for lowly animal needs into a complete rejection of the earthly city, which perhaps was becoming a burden for the late Romans. Thereafter you had Guenon's division of spiritual authority and temporal power, with only the spiritual being truly public. But even in the temporal world the proper distinction between public and private persisted and was widely understood.

Arendt identifies some modern inversals that even Evola did not seem to catch. She says that the modern devolution of power actually destroyed the idea of ``public good'' entirely in the West and replaced it with a confluence of private interests which slowly corrupts and devalues our sense of the public sphere. Furthermore, she identifies this counter-tradition with ``society''! As opposed to the classical Senators and medieval peers, whose speech was truly free, modernity more and more requires people to agree with each other to achieve social graces. ``The victory of equality in the modern world is only the political and legal recognition of the fact that society has conquered the public realm, and that distinction and difference have become private matters of the individual.'' (41)

The neopagan-Christian distinction is therefore obvious. Modernity basically tells religion and all other former public interests to step aside and let the private interests rule. Christians earnestly believe that they can accomplish their goals while following the rules of society. The Nouvelle Droite, from a utilitarian perspective at least, sees the Christian ideal as a waste of time, and wants nothing to do with ``society'', because a true nation needs a national spirit. Yet I still do not believe that the neopagans are right from the deeper, spiritual perspective.

By the way, I recently read a Japanese spiritualist's idea of Christianity where he focused on the concept of ``not peace, but a sword'' and claimed that Christians should also be enemies to this world. It's certainly all there in the Bible, but it would be a very Protestant and un-traditional rereading.

\hfill

\texttt{Jason-adam on 2013-04-12 at 14:40 said:}

I am glad to see my question brought out good points by both Cologero and Avery... ... I dont recall Evola discussing ever the concept of society but Avery's comments reminded me of Carl Schmitt discussing the differences, tension of, and basic opposition to state and society. The modern world is clearly the same thing as the liberal world and liberalism believes that the central place once occupied by spirituality should be replaced by private inquiry and skepticism. It is plain to see how this means the revolt against tradition. The catholic Church fought against this idea and tried to subjugate society to itself but after Vatican II it has accepted a role as punching bad of society and a refuge for old women nd closeted faggots... ... of course Schmitt was modern himself in believing that the state could be a substitute for spiritual authority which was the main flaw in fascism. i'm curious to see how Gentile describes the situation.

I am not a neopagan because, besides my feeling that catholicism explains the world in a better way than paganism can, my parents are catholics, my grandparents are catholics and my family has a heritage of being catholic knights. this may sound stupid but i dont feel a need to rebel against my family's history, even if the entire Church visible has become apostate, i dont see a need to abandon my belifs. if need be i can live my life by myself separate from the world in the truest christian sense of the word.

This actually raises another question ; before my abandoning national bolshevism, I was constantly attacked by my former `comrades' as being sentimental and reactionary for my views -- so would you say that attachment to family and ancestors is emotional and should be overcome ? or is devotion to family the essence of Tradition ?

\hfill

\texttt{Cologero on 2013-04-12 at 17:41 said:}

I'm not grokking this Avery. For example, just from the first paragraph. Evola criticized Christianity for sacralizing labour; he himself avoided it on the principles you noted. Did Christian really invent ``disgust for lowly animal needs''? Are you making defecation a pagan virtue? Are you saying that Guenon somehow invented or fabricated the concepts of spiritual authority and temporal power? There is no ``thereafter'', it always was. Where does Guenon imply that only the spiritual authority is truly public?

Arendt adds nothing to this discussion and actually seems to be harming your intellectuality. A ``confluence of private interests'' is a professorial expression for unbridled individualism, or a ``heap'' rather than an organic society. Evola has definitely ``caught'' that, in ways that Arendt and her admirers cannot understand.

Christians ``earnestly believe'' they should be following the will of God, not the rules of society.

You are partly right about the new right. There is no ``nation'' without a national spirit; in that case, it is just a country. They reject Christianity as the expression of that spirit. It is an ideological position, not a utilitarian one.

\hfill

\texttt{Michael on 2013-04-12 at 19:58 said:}

Just curious, where does identifying Christians with brahmins and neo-pagans with kshatriyas come from?

\hfill

\texttt{Avery Morrow on 2013-04-12 at 21:31 said:}

\emph{Did Christian really invent ``disgust for lowly animal needs''?}

No, you've got it backwards. The pagans saw politics as part of the intellectual life. The Christians didn't. That's all I meant to say.

Arendt doesn't stick to Evola's line but I'm pretty sure she has got more figured out than you give her credit for: ``The modern reversal followed and left unchallenged the most important reversal with which Christianity had broken into the ancient world ... The victory of the Christian faith in the ancient world was largely due to this reversal of the individual life for the body politic , which brought hope to those who knew that their world was doomed ... This reversal could not but be disastrous for the esteem and the dignity of politics.'' (314)

\hfill

\texttt{Cologero on 2013-04-13 at 16:45 said:}

It was just being used metaphorically, Michael, to show how one-dimensional the Christians had become.

\hfill

\end{sffamily}
\end{footnotesize}

